\section{The Cloud-Native Observability Maturity Model (CN-OMM)}
\label{sec:cn_omm}

To overcome the theoretical and architectural deficiencies identified in Section~\ref{sec:critical_analysis}, we formulate the \textbf{Cloud-Native Observability Maturity Model (CN-OMM)}. Following De Bruin et al.'s top-down hierarchical architecture (\textit{Domain $\rightarrow$ Orthogonal Dimensions $\rightarrow$ Sub-Components}) \cite{debruin_2005} and Lasrado et al.'s structural canon \cite{lasrado_2015}, this section presents: (i) the mathematical formalization of cloud-native observability and its dialectical constraints, (ii) the five orthogonal capability dimensions ($D_1$--$D_5$) and their 15 sub-components, and (iii) the five evolutionary maturity stages and their configurational transition rules. Figure~\ref{fig:cn_omm_architecture} illustrates the complete methodological and structural architecture of the CN-OMM.

\begin{figure*}[t]
\centering
\begin{tikzpicture}[
    font=\sffamily\scriptsize,
    >=Latex,
    box/.style={draw=black!75, rounded corners=3pt, align=center, minimum height=0.85cm, inner sep=4pt, line width=0.6pt},
    dimbox/.style={box, fill=blue!6, minimum width=2.85cm, minimum height=1.05cm},
    lvlbox/.style={box, fill=teal!8, minimum width=2.85cm, minimum height=0.95cm},
    gatebox/.style={box, fill=amber!12, draw=orange!80!black, minimum width=3.1cm, minimum height=1.15cm},
    evalbox/.style={box, fill=purple!7, draw=purple!70!black, minimum width=2.85cm, minimum height=1.05cm},
    arr/.style={->, thick, draw=black!75}
]

% Column 1: 5 Orthogonal Dimensions
\node[font=\sffamily\bfseries\footnotesize, anchor=south] at (0, 3.0) {1. Five Orthogonal Dimensions};
\node[dimbox] (d1) at (0, 2.2) {\textbf{$D_1$: Telemetry Depth}\\($M, L, T, E, P$ + eBPF)};
\node[dimbox] (d2) at (0, 1.0) {\textbf{$D_2$: Semantic Context}\\(W3C + Exemplars + $T_{\text{self}}$)};
\node[dimbox] (d3) at (0, -0.2) {\textbf{$D_3$: Diagnostic RCA}\\($MTTDiag, Acc$, Bias Free)};
\node[dimbox] (d4) at (0, -1.4) {\textbf{$D_4$: Resilience \& SLOs}\\(Chaos Probes \& Budgets)};
\node[dimbox] (d5) at (0, -2.6) {\textbf{$D_5$: Probe Economics}\\($\Delta_{\text{CPU}} < 15\%$, Index-Free)};

% Column 2: Automated Assessment Engine
\node[font=\sffamily\bfseries\footnotesize, anchor=south] at (4.5, 3.0) {2. Automated Assessment};
\node[gatebox] (qci) at (4.5, 1.45) {\textbf{15 Declarative QCIs}\\PromQL, OTLP Schema,\\K8s \& Chaos CRD Audit};
\node[gatebox] (dialectic) at (4.5, -0.2) {\textbf{Dialectical Constraints}\\Diagnostic Info $I(\mathbf{x};\mathcal{O})$\\vs.\ Probe Overhead $\Omega(\mathcal{O})$};
\node[gatebox] (stagegate) at (4.5, -1.85) {\textbf{Non-Compensatory Gate}\\$\text{Level}_{\text{eff}} = \min_k \lfloor C(D_k) \rfloor$\\+ Euclidean Index $\mathcal{I}(\mathcal{S})$};

% Column 3: 5 Evolutionary Maturity Stages
\node[font=\sffamily\bfseries\footnotesize, anchor=south] at (9.0, 3.0) {3. Evolutionary Stage-Gates};
\node[lvlbox, fill=gray!12] (l1) at (9.0, 2.2) {\textbf{Level 1: Opaque}\\($\mathcal{O} = \emptyset$, Post-Mortem)};
\node[lvlbox, fill=red!8] (l2) at (9.0, 1.0) {\textbf{Level 2: Reactive CLI}\\(\texttt{kubectl}, $Acc = 20\%$)};
\node[lvlbox, fill=orange!10] (l3) at (9.0, -0.2) {\textbf{Level 3: Fragmented}\\(Siloed Tracing, $Acc = 0\%$)};
\node[lvlbox, fill=green!14, draw=green!50!black, line width=1.0pt] (l4) at (9.0, -1.4) {\textbf{Level 4: Correlated}\\(PANOPTES, $Acc = 83.3\%$)};
\node[lvlbox, fill=blue!12] (l5) at (9.0, -2.6) {\textbf{Level 5: Autonomous}\\(Adaptive AIOps Loop)};

% Column 4: Multi-Method DSR Validation
\node[font=\sffamily\bfseries\footnotesize, anchor=south] at (13.3, 3.0) {4. DSR \& Empirical Validation};
\node[evalbox] (v1) at (13.3, 1.45) {\textbf{Helgesson Type-1 Eval.}\\Becker R1--R8 Audit vs.\\7 Academic/Vendor Models};
\node[evalbox] (v2) at (13.3, -0.2) {\textbf{Helgesson Type-3 Eval.}\\$N=30$ Paired Chaos Trials\\($MTTDiag: 300\text{s} \rightarrow 17.3\text{s}$)};
\node[evalbox] (v3) at (13.3, -1.85) {\textbf{Probe Effect Verification}\\$N=60$ Injector Trials\\($\Delta_{\text{CPU}} = 9.51\% < 15\%$)};

% Connections from D1-D5 to Assessment
\draw[arr] (d1.east) -- (qci.west);
\draw[arr] (d2.east) -- (qci.west);
\draw[arr] (d2.east) -- (dialectic.west);
\draw[arr] (d3.east) -- (dialectic.west);
\draw[arr] (d4.east) -- (stagegate.west);
\draw[arr] (d5.east) -- (dialectic.west);
\draw[arr] (d5.east) -- (stagegate.west);

\draw[arr] (qci.south) -- (dialectic.north);
\draw[arr] (dialectic.south) -- (stagegate.north);

% Connections from Assessment to Levels
\draw[arr] (stagegate.east) -- (l2.west);
\draw[arr] (stagegate.east) -- (l3.west);
\draw[arr] (stagegate.east) -- (l4.west);
\draw[arr] (l1.south) -- (l2.north);
\draw[arr] (l2.south) -- (l3.north);
\draw[arr] (l3.south) -- (l4.north);
\draw[arr] (l4.south) -- (l5.north);

% Connections to Validation
\draw[arr] (l2.east) -- (v2.west);
\draw[arr] (l3.east) -- (v2.west);
\draw[arr] (l4.east) -- (v2.west);
\draw[arr] (l4.east) -- (v3.west);
\draw[arr] (qci.east) to[out=15, in=165] (v1.west);

\end{tikzpicture}
\caption{Methodological and Architectural Overview of the Cloud-Native Observability Maturity Model (CN-OMM), Integrating Five Orthogonal Dimensions ($D_1$--$D_5$), Automated Declarative Auditing (15 QCIs), Non-Compensatory Stage-Gate Scoring, and Type-1/Type-3 Empirical Chaos Validation.}
\label{fig:cn_omm_architecture}
\end{figure*}


\subsection{Mathematical Formalization of Distributed State Observability}
\label{subsec:math_formalization}

Let a cloud-native microservice system orchestrated on Kubernetes be modeled over discrete observation windows $t \in \mathbb{T}$ as a dynamic directed call graph $\mathcal{G}(t) = (\mathcal{V}(t), \mathcal{E}(t))$, where $\mathcal{V}(t) = \{v_1, v_2, \dots, v_K\}$ denotes the set of active service workloads (Pods/containers) and $\mathcal{E}(t) \subseteq \mathcal{V}(t) \times \mathcal{V}(t)$ denotes synchronous or asynchronous invocation edges. Each workload $v_k \in \mathcal{V}(t)$ possesses a latent internal operational state $\mathbf{x}_k(t) \in \mathcal{X}_k$ (capturing CPU cgroup throttling, memory heap pressure, socket queue saturation, and application exception states), forming the global system state $\mathbf{x}(t) = [\mathbf{x}_1(t), \dots, \mathbf{x}_K(t)]^\top$.

The external telemetry emitted by the system during window $t$ is defined as the multimodal five-tuple:
\begin{equation}
    \mathcal{O}(t) = \langle M(t), L(t), T(t), E(t), P(t) \rangle
    \label{eq:telemetry_tuple}
\end{equation}
where $M(t)$ represents aggregated numerical time-series metrics, $L(t)$ represents discrete structured or unstructured log records, $T(t)$ represents distributed execution traces composed of causal spans $s_i = (\text{trace\_id}, \text{span\_id}, \text{parent\_id}, v_k, t_{\text{start}}, t_{\text{end}})$, $E(t)$ represents Kubernetes control-plane and scheduler state transitions, and $P(t)$ represents continuous kernel/runtime execution profiles (e.g., eBPF stack samples).

\subsubsection{The Dialectical Trade-Off: Diagnostic Information vs.\ Probe Effect}
Given a fault injection or production incident where root cause $v^* \in \mathcal{V}(t)$ experiences an anomalous state transition, a diagnostic operator $\Phi: \mathcal{O}(t) \rightarrow \mathcal{V}(t) \cup \{\varnothing\}$ attempts to infer $\hat{v} = \Phi(\mathcal{O}(t))$ within diagnostic latency $MTTDiag$. In information-theoretic terms, the diagnostic utility of an observability architecture is proportional to the mutual information $I(\mathbf{x}(t); \mathcal{O}(t)) = H(\mathbf{x}(t)) - H(\mathbf{x}(t) \mid \mathcal{O}(t))$, which measures the reduction in diagnostic uncertainty $H(\mathbf{x} \mid \mathcal{O})$.

However, acquiring and transporting $\mathcal{O}(t)$ incurs a physical \textbf{Probe Effect Vector} $\Omega(\mathcal{O}(t))$:
\begin{equation}
    \Omega(\mathcal{O}(t)) = \left\langle \Delta_{\text{CPU}}(\mathcal{O}), \; \Delta_{\text{RAM}}(\mathcal{O}), \; \Delta_{\text{Net}}(\mathcal{O}), \; \Delta_{\text{Storage}}(\mathcal{O}) \right\rangle
    \label{eq:probe_effect}
\end{equation}
where $\Delta_{\text{CPU}}(\mathcal{O}) = (\text{CPU}_{\text{inst}} - \text{CPU}_{\text{base}}) / \text{CPU}_{\text{alloc}}$ denotes the marginal fractional CPU overhead imposed by SDKs, eBPF probes, and OpenTelemetry collectors. Maturity progression is therefore governed by a constrained optimization problem: maximizing diagnostic reconstructibility $I(\mathbf{x}(t); \mathcal{O}(t))$ while strictly satisfying the stability envelope $\Omega(\mathcal{O}(t)) \preceq \boldsymbol{\theta}_{\max}$ (e.g., $\Delta_{\text{CPU}} < \theta_{\text{CPU}} = 15\%$).

\subsubsection{Formalizing Cascading Span Bias and Exclusive Self-Time ($T_{\text{self}}$)}
A critical theoretical insight of the CN-OMM is explaining why monomodal distributed tracing fails under synchronous call chains. For any span $s_i \in T(t)$ executing on service $v_k$ with synchronous child spans $\text{children}(s_i) = \{s_j \in T(t) \mid \text{parent\_id}(s_j) = \text{span\_id}(s_i)\}$, its inclusive duration is:
\begin{equation}
    D(s_i) = t_{\text{end}}(s_i) - t_{\text{start}}(s_i) = T_{\text{self}}(s_i) + \sum_{s_j \in \text{children}(s_i)} D(s_j)
    \label{eq:inclusive_duration}
\end{equation}
Whenever a leaf or intermediate service $v^*$ suffers resource throttling or network latency ($\Delta \tau > 0$), every upstream ancestor $s_{\text{anc}}$ along the path from the ingress gateway ($v_{\text{root}}$) to $v^*$ exhibits an inflated inclusive duration $D(s_{\text{anc}}) \ge D(s^*) + \Delta \tau$. Consequently, any monomodal diagnostic operator $\Phi_{\text{trace}}$ that ranks suspects by inclusive tail latency ($P_{95}/P_{99}$ of $D(s_i)$) systematically misattributes the fault to upstream callers ($\hat{v} \neq v^*$), yielding \textbf{Cascading Span Bias}. Eliminating this bias requires decomposing inclusive duration into \textbf{exclusive span self-time}:
\begin{equation}
    T_{\text{self}}(s_i) = \max\left(0, \; D(s_i) - \sum_{s_j \in \text{children}(s_i)} \mathcal{U}(s_j)\right)
    \label{eq:exclusive_selftime}
\end{equation}
where $\bigcup_{s_j} \mathcal{U}(s_j)$ denotes the union of non-overlapping child execution intervals, coupled via shared W3C \texttt{trace\_id} keys to co-occurring log errors $L(t)$ and container resource metrics $M(t)$.

\subsection{The Five Orthogonal Dimensions and Fifteen Sub-Components}
\label{subsec:dimensions}

Consistent with Miller's cognitive capacity principle ($7 \pm 2$ categories) emphasized by Lasrado et al.\ \cite{lasrado_2015} and De Bruin et al.'s component layering \cite{debruin_2005}, the CN-OMM structures cloud-native observability across \textbf{five orthogonal dimensions ($D_1$--$D_5$)}, each decomposed into three measurable sub-components ($d_{k.1}, d_{k.2}, d_{k.3}$), totaling 15 architectural capabilities (summarized in Table~\ref{tab:cn_omm_dimensions}):
\begin{itemize}[leftmargin=*, nosep]
    \item \textbf{$D_1$ -- Telemetry Acquisition \& Instrumentation Depth:} Evaluates the breadth of emitted signal modalities ($d_{1.1}$: $M, L, T, E, P$), adherence to open, vendor-neutral collection protocols ($d_{1.2}$: OpenTelemetry OTLP and W3C TraceContext), and visibility into kernel-level socket and scheduler operations via non-intrusive eBPF probes ($d_{1.3}$).
    \item \textbf{$D_2$ -- Semantic Context \& Cross-Pillar Correlation:} Assesses the structural coupling across telemetry streams, encompassing automatic Kubernetes topological enrichment ($d_{2.1}$: Pod UID, Node, Deployment labels), bi-directional context propagation ($d_{2.2}$: W3C \texttt{trace\_id} injection into structured JSON logs and OpenMetrics Exemplars linking histograms to traces), and causal self-time graph isolation ($d_{2.3}$: $T_{\text{self}}$ computation).
    \item \textbf{$D_3$ -- Diagnostic \& Analytical Intelligence (RCA):} Captures the operational efficacy of fault management, measured through automated anomaly detection latency ($d_{3.1}$: $MTTD$), multi-pillar root-cause localization speed and accuracy ($d_{3.2}$: $MTTDiag, Acc$), and immunity to cascading upstream misattribution alongside explainable diagnostic evidence ($d_{3.3}$).
    \item \textbf{$D_4$ -- Resilience Verification \& SLO Governance:} Bridges observability with Reliability and Chaos Engineering \cite{basiri_2016, beyer_2016}, tracking multi-window SLI/SLO error-budget enforcement ($d_{4.1}$), continuous verification of telemetry pipelines via declarative Chaos Engineering steady-state probes ($d_{4.2}$), and closed-loop automated remediation triggers ($d_{4.3}$).
    \item \textbf{$D_5$ -- Instrumentation Economics \& Probe Effect Governance:} Formalizes operational sustainability and FinOps discipline, auditing steady-state CPU and memory overhead budgets ($d_{5.1}$: $\Delta_{\text{CPU}}, \Delta_{\text{RAM}} < 15\%$), intelligent tail-based or anomaly-aware sampling ($d_{5.2}$), and index-free object storage architectures that prevent high-cardinality index explosion ($d_{5.3}$).
\end{itemize}

\begin{table*}[t]
\centering
\caption{The Five Orthogonal Dimensions ($D_1$--$D_5$), Fifteen Sub-Components ($d_{k.j}$), and Automated Quantitative Capability Indicators (QCIs) of the CN-OMM}
\label{tab:cn_omm_dimensions}
\footnotesize
\setlength{\tabcolsep}{4pt}
\renewcommand{\arraystretch}{1.18}
\begin{tabularx}{\textwidth}{l l >{\raggedright\arraybackslash}p{4.3cm} >{\raggedright\arraybackslash}X l}
\toprule
\textbf{Dim.} & \textbf{Code} & \textbf{Sub-Component Name} & \textbf{Quantitative Capability Indicator (QCI) \& Declarative Cluster Audit Rule} & \textbf{Target Metric / Threshold} \\
\midrule
\multirow{3}{*}{\textbf{$D_1$}} & $d_{1.1}$ & Multi-Signal Telemetry Coverage & Fraction of active Kubernetes workloads emitting structured $M, L, T, E,$ and continuous profiles $P$. & $\text{QCI}_{1.1} = |\mathcal{V}_{\text{inst}}| / |\mathcal{V}_{\text{total}}| \ge 0.95$ \\
 & $d_{1.2}$ & Vendor-Neutral OTel Standardization & Proportion of telemetry pipelines using open OTLP/W3C schemas without proprietary lock-in. & $\text{QCI}_{1.2} = \text{OTLP}_{\text{flows}} / \text{Total}_{\text{flows}} = 1.0$ \\
 & $d_{1.3}$ & Kernel \& Runtime Depth (eBPF) & Coverage of non-intrusive kernel/socket/cgroup telemetry (e.g., Beyla/Cilium eBPF probes). & $\text{QCI}_{1.3} \in [0, 1]$ (L4 $\ge 0.80$) \\
\midrule
\multirow{3}{*}{\textbf{$D_2$}} & $d_{2.1}$ & Topological Metadata Enrichment & Automatic injection of K8s resource attributes (\texttt{k8s.pod.uid}, \texttt{node}, \texttt{deployment}) across signals. & $\text{QCI}_{2.1} = \text{Enriched}_{\text{records}} / \text{Total} \ge 0.99$ \\
 & $d_{2.2}$ & Cross-Pillar Context Coupling & Rate of structured logs containing valid W3C \texttt{trace\_id}/\texttt{span\_id} and metrics exposing Exemplars. & $\text{QCI}_{2.2} = \min(\rho_{\text{log-trace}}, \rho_{\text{exemplar}}) \ge 0.90$ \\
 & $d_{2.3}$ & Causal Self-Time Isolation ($T_{\text{self}}$) & Capability to decompose inclusive span duration $D(s_i)$ into exclusive execution $T_{\text{self}}(s_i)$. & $\text{QCI}_{2.3} \in \{0, 0.5, 1.0\}$ (L4 $= 1.0$) \\
\midrule
\multirow{3}{*}{\textbf{$D_3$}} & $d_{3.1}$ & Automated Anomaly DetectionLatency & Mean/median time to detect steady-state deviation ($MTTD$) without manual CLI polling. & $\text{QCI}_{3.1}: MTTD \le 15.0\,\text{s}$ \\
 & $d_{3.2}$ & Root-Cause Localization ($MTTDiag, Acc$) & Automated multi-pillar fault localization accuracy ($Acc$) and diagnostic latency ($MTTDiag$). & $\text{QCI}_{3.2}: Acc \ge 80\%, MTTDiag \le 30\,\text{s}$ \\
 & $d_{3.3}$ & Cascading Bias Immunity \& Explainability & Resistance to upstream cascading span misattribution ($1 - \text{FPR}_{\text{cascade}}$) and auditability of RCA path. & $\text{QCI}_{3.3}: \text{FPR}_{\text{cascade}} \le 0.15$ \\
\midrule
\multirow{3}{*}{\textbf{$D_4$}} & $d_{4.1}$ & SLI/SLO \& Error Budget Governance & Percentage of critical user journeys governed by multi-window burn-rate SLOs. & $\text{QCI}_{4.1} = |\text{Services}_{\text{SLO}}| / |\mathcal{V}_{\text{crit}}| \ge 0.90$ \\
 & $d_{4.2}$ & Active Chaos Engineering Verification & Continuous execution of declarative fault injection with automated PromQL/HTTP steady-state probes. & $\text{QCI}_{4.2}: \text{ChaosProbes}_{\text{active}} \ge 1.0$ \\
 & $d_{4.3}$ & Closed-Loop Remediation Readiness & Integration of observability alerts/RCA with automated rollback, scaling, or traffic shifting gates. & $\text{QCI}_{4.3} \in [0, 1]$ (L5 $\ge 0.90$) \\
\midrule
\multirow{3}{*}{\textbf{$D_5$}} & $d_{5.1}$ & CPU \& Memory Probe-Effect Budget & Marginal steady-state resource overhead ($\Delta_{\text{CPU}}, \Delta_{\text{RAM}}$) imposed by collectors and SDKs. & $\text{QCI}_{5.1}: \Delta_{\text{CPU}} < 15\%, \Delta_{\text{RAM}} < 15\%$ \\
 & $d_{5.2}$ & Adaptive \& Tail-Based Sampling & Ratio of anomalous/high-latency traces retained vs.\ redundant baseline telemetry dropped. & $\text{QCI}_{5.2}: \text{Recall}_{\text{anomaly}} \ge 0.99$ \\
 & $d_{5.3}$ & Index-Free Storage \& Cardinality Control & Storage amplification factor and resistance to high-cardinality label explosion (e.g., Loki/Tempo). & $\text{QCI}_{5.3}: \text{IndexRatio} \le 0.05$ \\
\bottomrule
\end{tabularx}
\end{table*}


\subsection{The Five Evolutionary Maturity Levels and Stage-Gate Matrix}
\label{subsec:maturity_levels}

Intersecting the five orthogonal dimensions with five evolutionary stages yields the $5 \times 5$ CN-OMM Architectural Matrix presented in Table~\ref{tab:cn_omm_matrix}. Each level represents a distinct architectural regime characterized by formal state properties:

\begin{table*}[t]
\centering
\caption{The Cloud-Native Observability Maturity Model (CN-OMM) $5\text{-Level} \times 5\text{-Dimension}$ Architectural Matrix}
\label{tab:cn_omm_matrix}
\footnotesize
\setlength{\tabcolsep}{3.5pt}
\renewcommand{\arraystretch}{1.22}
\begin{tabularx}{\textwidth}{l >{\raggedright\arraybackslash}X >{\raggedright\arraybackslash}X >{\raggedright\arraybackslash}X >{\raggedright\arraybackslash}X >{\raggedright\arraybackslash}X}
\toprule
\textbf{Maturity Level} & \textbf{$D_1$: Telemetry Acquisition \& Depth} & \textbf{$D_2$: Semantic Context \& Correlation} & \textbf{$D_3$: Diagnostic \& Analytical Intelligence} & \textbf{$D_4$: Resilience Verification \& SLOs} & \textbf{$D_5$: Instrumentation Economics ($\Omega$)} \\
\midrule
\textbf{Level 1: Opaque} \newline \textit{(Blind Execution)} & $\mathcal{O} = \emptyset$ or unstructured \texttt{stdout} ephemeral container streams. & Zero metadata enrichment; telemetry lost upon Pod eviction or restart. & Post-mortem human inspection triggered by external user complaints. & No SLIs/SLOs; absence of health probes or resilience testing. & Unmeasured ($0\%$ intentional telemetry, high outage cost). \\
\midrule
\textbf{Level 2: Reactive} \newline \textit{(Disjoint Polling)} & Infrastructure polling (\texttt{kubelet}/cAdvisor) and ad-hoc CLI (\texttt{kubectl top/logs}). & Disjoint local streams; no shared timestamps or request identifiers. & Manual symptom triage; high cognitive load ($MTTDiag \ge 300\,\text{s}$, $Acc \le 20\%$). & Static host/Pod threshold alerts; high alert fatigue; manual chaos tests. & Ad-hoc log retention; uncontrolled burst overhead during CLI scraping. \\
\midrule
\textbf{Level 3: Fragmented} \newline \textit{(Siloed Pillars)} & Centralized $M, L, T$ via agents/SDKs, but lacking kernel/eBPF visibility. & Siloed databases (Prometheus, Elasticsearch, Jaeger) queried via manual time-window pivoting. & Monomodal trace/metric sorting; vulnerable to \textbf{Cascading Span Bias} ($Acc \rightarrow 0\%$ on resource faults). & Service-level SLIs defined, but chaos experiments lack automated telemetry gates. & Full-index storage bloat; static head-based sampling drops tail anomalies; uncalibrated $\Delta_{\text{CPU}}$. \\
\midrule
\textbf{Level 4: Correlated} \newline \textit{(Multimodal Fusion)} \newline \textbf{[PANOPTES Target]} & Unified OTel + eBPF zero-code collectors capturing $M, L, T, E, P$ with $100\%$ OTLP compliance. & Mandatory W3C TraceContext in logs, Prometheus Exemplars, and exclusive self-time ($T_{\text{self}}$) graph isolation. & Deterministic cross-pillar RCA ($Acc = 83.3\%$, $MTTDiag = 17.3\,\text{s}$); immune to cascading trace propagation bias. & Error-budget SLOs continuously verified via declarative Chaos Mesh / Litmus steady-state probes. & Strictly bounded Probe Effect ($\Delta_{\text{CPU}} = 9.51\% < 15\%$); index-free storage (Loki/Tempo) and tail sampling. \\
\midrule
\textbf{Level 5: Autonomous} \newline \textit{(Adaptive AIOps)} & Context-aware dynamic kernel/runtime instrumentation modulated in real time. & Unified topological lakehouse with continuous causal graph discovery. & Explainable GNN/LLM root-cause reasoning with counterfactual confidence bounds. & Autonomous self-healing (canary rollback, cgroup scaling, circuit breaking) driven by chaos verification. & FinOps/GreenOps dynamic signal weighting ($w_i = \alpha_i S_i$) optimizing diagnostic entropy per CPU watt. \\
\bottomrule
\end{tabularx}
\end{table*}


\subsubsection{Level 1: Opaque (Blind Execution)}
At Level 1, the system operates as a black box ($\mathcal{O}(t) \approx \emptyset$). Workloads emit unstructured text to local container \texttt{stdout} files that vanish upon Pod eviction or node rescheduling. Because $\text{rank}(\mathcal{O}) = 0$, internal fault states cannot be reconstructed; outages are discovered externally by end users, and post-mortem debugging relies on manual reproduction attempts.

\subsubsection{Level 2: Reactive Monitoring (Disjoint Infrastructure Polling)}
At Level 2, platform teams collect infrastructure-level utilization counters (\texttt{kubelet}/cAdvisor CPU and memory metrics) and perform reactive manual triage via Kubernetes CLI utilities (\texttt{kubectl top}, \texttt{kubectl logs}, \texttt{kubectl describe pod}). While binary container failures (e.g., \texttt{OOMKilled} or \texttt{CrashLoopBackOff}) are visible through control-plane events $E(t)$, the architecture lacks request-scoped context or historical correlation. Consequently, under non-destructive degradations (such as network packet delay or cgroup CPU throttling), engineers face severe cognitive overload, resulting in timeouts ($MTTDiag \ge 300\,\text{s}$) and low localization accuracy.

\subsubsection{Level 3: Fragmented Pillar Observability (Siloed Signals)}
At Level 3, organizations deploy dedicated collection agents for metrics (e.g., Prometheus), logs (e.g., Elasticsearch/Fluentd), and distributed traces (e.g., Jaeger/Zipkin). Despite high financial investment and storage footprint, these signals reside in \textbf{isolated analytical silos} ($D_2 < 4$). Logs lack embedded W3C \texttt{trace\_id} fields, metrics lack Exemplars, and trace engines evaluate spans using raw inclusive duration $D(s_i)$ rather than $T_{\text{self}}(s_i)$. As proven in Section~\ref{sec:empirical_validation}, Level 3 systems fall victim to the \textit{Monomodal Additivity Fallacy}: under synchronous microservice cascades, monomodal trace sorting achieves $0\%$ root-cause accuracy due to \textit{Cascading Span Bias}, while full-text log indexing inflates storage and CPU overhead ($D_5$).

\subsubsection{Level 4: Context-Correlated Multimodal Observability (Semantic Fusion)}
Level 4 represents the deterministic synthesis of the dialectical tension between visibility and overhead (instantiated by our \textbf{PANOPTES} reference architecture). All signals are unified under OpenTelemetry (OTLP) and enriched with zero-code eBPF kernel telemetry ($D_1 \ge 4$). Crucially, Level 4 enforces strict cross-pillar context coupling ($D_2 \ge 4$): every structured log entry carries active W3C \texttt{trace\_id}/\texttt{span\_id} attributes, metric histograms expose OpenMetrics Exemplars, and the diagnostic engine computes exclusive span self-time $T_{\text{self}}(s_i)$ fused with container resource saturation and log error rates. Furthermore, observability pipelines are actively verified by declarative Chaos Engineering steady-state probes ($D_4 \ge 4$) and backed by index-free storage (Grafana Loki and Tempo) that constrains steady-state CPU overhead below $\theta_{\text{CPU}} = 15\%$ ($D_5 \ge 4$).

\subsubsection{Level 5: Autonomous \& Adaptive Self-Healing (Closed-Loop AIOps)}
At Level 5, the correlated multimodal foundation of Level 4 is extended into a continuous, closed-loop cybernetic control system. Telemetry sampling and signal weights ($w_i = \alpha_i S_i$ \cite{amgothu_2025}) are dynamically modulated in real time based on error-budget burn rates and FinOps/GreenOps carbon constraints ($D_5 = 5$). Root-cause localization leverages explainable Graph Neural Networks (GNNs) or neuro-symbolic LLM agents grounded in the causal $T_{\text{self}}$ topology ($D_3 = 5$), automatically triggering verified Kubernetes remediation workflows (canary rollbacks, dynamic cgroup vertical scaling, or service-mesh circuit breaking) without human intervention ($D_4 = 5$).

\subsection{Developmental Lifecycle and Empirical Population (De Bruin et al., 2005)}
\label{subsec:debruin_lifecycle}

To ensure methodological rigor and repeatability, the construction of the CN-OMM operationalizes the six sequential lifecycle phases established by De Bruin et al.\ \cite{debruin_2005}:
\begin{itemize}[leftmargin=*, nosep]
    \item \textbf{1. Scope:} The model's domain scope is formally bounded to cloud-native, polyglot microservice architectures orchestrated by Kubernetes, addressing the architectural challenges of multi-signal telemetry ingestion, cross-pillar semantic correlation, and instrumentation overhead governance.
    \item \textbf{2. Design:} We adopted a top-down, multi-dimensional stage-gate architecture structured across five orthogonal capability dimensions ($D_1$--$D_5$), enforcing non-compensatory maturation criteria rather than additive arithmetic averages.
    \item \textbf{3. Populate:} The capability sub-components ($d_{k.j}$) were populated via a hybrid approach: top-down control-theoretic reconstructibility axioms (Kalman \cite{kalman_1960}) combined with bottom-up empirical evidence synthesized from a Rapid Review of the cloud-native literature \cite{cartaxo_2018}. Out of 312 initial records retrieved across IEEE Xplore, ACM Digital Library, SpringerLink, and ScienceDirect, 10 primary studies on Kubernetes observability were selected (\cite{kosinska_2023, usman_2022, faseeha_2025, araujo_2024, hartikainen_2024, gomes_2025, han_2024, deng_2024, zanatta_23, wei_2025}). These studies directly grounded the definition of essential components ($D_1$), correlation gaps ($D_2$), diagnostic bottlenecks ($D_3$), and probe overhead constraints ($D_5$).
    \item \textbf{4. Test:} To overcome the pervasive ``Validation Gap'' documented by Wendler \cite{wendler_2012}, the CN-OMM undergoes dual validation (Helgesson et al.\ \cite{helgesson_2012}): a Type-1 structural and DSR requirement audit (Section~\ref{subsec:type1_evaluation}) and a Type-3 real-system empirical experiment under controlled Chaos Engineering (Section~\ref{subsec:empirical_results}).
    \item \textbf{5. Deploy:} The model is operationalized through executable Kubernetes auditing artifacts (Quantitative Capability Indicators, Section~\ref{sec:assessment}) and the open-source PANOPTES reference implementation, enabling automated cluster maturity grading.
    \item \textbf{6. Maintain:} The schema is designed for continuous evolution, where Level 4 provides the semantically linked, bias-free multimodal telemetry substrate necessary to train and operate Level 5 closed-loop AIOps and GreenOps models.
\end{itemize}
